\documentclass{reqengbook}

\title{
\bf\sffamily\fontsize{22}{28}\selectfont
Requirements Engineering for Software Developers
}

\author{\sffamily\fontsize{20}{30}\selectfont Björn Regnell}

\date{\vspace{2em}\sffamily\small Updated: \today 
\\ License: CC-BY-SA 
\\ \url{https://github.com/bjornregnell/reqeng-book} 
}

\begin{document}
%\pagenumbering{roman} 

\begin{minipage}{0.5\textwidth}
  \maketitle%
\end{minipage}%
\begin{minipage}{0.5\textwidth}
  \centering\vspace{2.2cm}
  \includegraphics[width=0.9\textwidth]{../img/phone-support.jpg}%
\end{minipage}%
\pagebreak

\setcounter{tocdepth}{2}

\begin{multicols*}{3}  % * means unbalanced columns allowed
\fontsize{8.4}{11}\selectfont

\tableofcontents
\end{multicols*}

\newgeometry{hmargin={45mm,148.5mm},vmargin={15mm,15mm}}

\part{Domain Knowledge}  %%%%%%%%%%%%%%%%%%%%%%%%%%%%

\chapter{Introduction}%
\MarginPage{%

\LectureOnly{\section{Introduction}}

\begin{Slide}{What is Requirements Engineering (RE)?}

\begin{itemize}
\item RE is focused on the
\begin{itemize}
\item \textbf{features} of software-intensive systems 
\item \textbf{system context}, including users and connected systems
\item \textbf{development context}, including stakeholders' intentions 

\end{itemize}
\item The RE process involves 
\begin{itemize}
\item knowledge-building \hfill research
\item consensus-building \hfill agree
\item decision-making    \hfill choose
\item innovation         \hfill generate ideas
\item communication      \hfill be pedagogical

\end{itemize}
\end{itemize}
\end{Slide}

\begin{Slide}{What is a requirement?}

\begin{itemize}
\item Something \textbf{needed} or \textbf{wanted}.

\item A documented \textbf{representation} of\\something needed or wanted.
\item The word ''requirement'' have different meanings:\\
  must, wish, idea, specific design, rationale, ...

\item The most \textbf{general meaning}:\\
  \textit{any} kind of \textbf{information entity} used in RE

\item You don't always get what you want and you often want things that you don't need...

\end{itemize}
\end{Slide}

\SwedishTermsStart
\SwedishTerm{requirements engineering}{kravhantering, kravteknik, kravanalys}
\SwedishTerm{requirement}{krav, önskemål}
\SwedishTerm{representation}{representation, avbildning}
\SwedishTerm{stakeholder}{intressent, aktör, sakägare}
}%
Requirements Engineering (RE) is a sub-discipline of Software Engineering (SE) that is focused on the \textit{potential features} of future software-intensive systems and their \textit{context}, which includes users and surrounding systems.

The requirements engineering process involves activities such as eliciting, specifying, validating, and selecting requirements. 
Requirements engineering is both a research discipline and an engineering practice, forming the foundation for our human ability to create useful software with high quality. When you work with anything related to the decision-making process of software development and its underlying intentions you are doing requirements engineering.

\section{What is a requirement?}%

The term 'requirement' often means \textbf{something needed or wanted}. It can also refer to a documented \emph{representation} of something needed or wanted. 

The term 'requirement' is a bit tricky as it is overloaded with several meanings. It is sometimes used as a very abstract term denoting \textit{any} type of information entity relevant to the intentions behind system development.  Sometimes we mean an optional wish or nice-to-have feature, sometimes we mean a mandatory feature without which the system would be pointless. Sometimes it refers to something explicitly decided and committed for delivery. Sometimes we mean the actual underlying need of a user that maybe still elusive and not explicitly stated, while we also may mean the carefully documented \textit{representation} of that underlying need. 

A documented representation of a requirement may very well unintendedly differ from the underlying actual need or wish, as our specification may be far from perfect. Hence, our views of an ''actual'' requirement and our views of its representation may be far from similar, if the representation is bad. It all depends on how humans interpret the representation.

Also, representations of requirements come in many different kinds and shapes at different levels of abstraction and detail, and hence, it is wise to be clear and specific about what kind of requirement we mean, to avoid that our usage of the term 'requirement' is misinterpreted. 


The subsequent sections aims to provide a concise overview of the set of terms that you can use when you think and communicate about different kinds of requirements in different requirements engineering contexts. This terminology is the starting-point for learning more about requirements engineering and developing skills needed to contribute to the invention of tomorrows' high-quality software that fulfill the human intentions behind its development.

\section{What is Requirements Engineering (RE)?}%
\subsection{Core activities of RE}%
\MarginPage{

\begin{Slide}{Main Activities of RE}

{\vspace{1em}\resizebox{!}{3em}{\hspace{0.5em}{\bf ?}\hspace{0.5em}\WritingHand\hspace{0.5em}\Checkedbox\hspace{0.5em}\LeftScissors}\vspace{0.5em}}

\begin{itemize}
\item The 4 main activities of RE: 
\begin{itemize}
\item \textbf{Elicitation} \hfill learning
\item \textbf{Specification} \hfill representing
\item \textbf{Validation}  \hfill checking
\item \textbf{Selection}   \hfill deciding
\end{itemize}
\item These activities are
\begin{itemize}
\item \textbf{Interdependent} \hfill output of one is input to other activities
\item \textbf{Concurrent} \hfill one activity you often trigger another
\item \textbf{Continuous} \hfill in the product life-cycle as software evolves

\end{itemize}
\end{itemize}
\end{Slide}
}% 
Working with requirements involves many different tasks and skills. The RE process is often characterized by these interrelated activities: 

\begin{itemize}
  \item \textbf{Elicitation}: to learn, invent and agree. We need to identify requirements before they can play a role in our software engineering process. Identifying requirements is far from trivial and we often need to work hard to discover the future needs of our users and sometimes we must be really creative and actually invent requirements representing our novel ideas on future software technology. Requirements elicitation is a learning and innovation process that provides the knowledge needed by other RE and SE activities. Elicitation also involves paving the way for agreement among stakeholders  if they have different views and conflicting goals, or at least elicit such conflicting goal as input to specification and selection.

  \item \textbf{Specification}: to represent, communicate and make persistent. It is not enough to just keep our knowledge about the evolving requirements in our heads. Requirements needs to be represented in a persistent way, so that we can retrieve and communicate them over time and space, as they evolve. The requirements models that are the results of requirements specification should be useful to all humans and tools involved in all dependent software engineering processes throughout the whole product life-cycle.
  \item \textbf{Validation}: to check that we have represented the right requirements in a good way. It is important to continuously validate that the current set of requirements are relevant and that their representation is of adequate quality. There is no such thing as a perfect and totally complete requirements specification, as that would require infinite resources and perfect knowledge about the future -- the major challenge is to tell when it is ''good enough'' for the involved stakeholders. With our requirements validation efforts we want to significantly reduce the risk of developing an unwanted system.
  \item \textbf{Selection}: to prioritize and decide what to include in the set of requirements that we make a commitment to implement, according to our product strategy and delivery plan, in relation to stakeholders' goals. Prioritization involves estimation using some measurement scale according to relevant selection criteria such as market value, development cost, requirements volatility, financial risk, etc. Delivery planning needs to take into account constraints on requirements implementation order, given that some requirements often depend on other requirements. 
\end{itemize} 

These activities are often conducted concurrently and in an iterative maner. When you specify or validate, you often discover missing requirements. When you elicit, you need to write down what you find in order not to forget all the things you learn. When you try to prioritize and make delivery plans, you may realize that you need to elicit more knowledge about requirements dependencies. Etcetera. Do not expect to carry out the above activities one at a time from top to bottom -- your RE work will be similar to a juggler having many balls in the air at the same time. 

\subsection{What is good RE?}

\begin{itemize}
  \item Cost-effective
  \item Happy stakeholders
  \item When are we ready?
\end{itemize}

\subsection{Requirements in development}
How are requirements specifications used in the development process?

\TODO Make a picture of the development process without assuming too much: perhaps cog wheels for each process and subprocess: Development(RE, design, code, test, integrate, deploy) and Operation 

\subsubsection{Requirements and product design}

\begin{itemize}
  \item Architecture: QR important basis for design choices
  \item Interaction design: based on contextual usage and ysability requirements
\end{itemize}


\subsubsection{Requirements and product verification}

\begin{itemize}
  \item Static and Dynamic testing (inspections and testing)
  \item Interface-based and structural testing (black box -- white box)
  \item Testing at different levels: unit -- system
  \item Regression testing: when evolving; verify that new or changed features do not brake existing features
  \item Acceptance testing: before delivery; verify fulfillment of user needs
\end{itemize}


\subsubsection{Requirements and product scoping}



\subsection{Requirements as constraints}

% \subsection{Decision quality}
% Or should this go into product scoping?

\section{Requirements Specification}



\subsection{What is a req spec?}


When we specify requirements we need to represent them in a useful way. We call such representations \textit{requirements models}. A requirement model is a \emph{simplified} representation of future versions of a system under development including its surrounding world -- we cannot expect to model every aspect with all its details, and we cannot fully predict every facet of all future stakeholders' intentions. We have limited time and resources so we need to make good-enough models that can reasonably support decision-making, design, verification, implementation and deployment in a cost-efficient way.

There are many aspects of requirements modeling and many different kinds of requirements. We need to model both functional and quality requirements, and model the requirements at different levels of abstraction and detail, while keeping in mind how the requirements models will be used by different stakeholders throughout a software product's whole life-cycle. Requirements are interrelated and can be part of a structure. Hence it is often useful to aggregate requirements into coherent bundles that each can be used as a single entity in requirements selection, while tracking its relations to other entities. Such interrelated requirements aggregates are often called \emph{features}. A set of features may be selected for implementation in a coming release based on priorities of stakeholders and available development resources. Therefore, we also need to represent priorities and resource constraints in our requirements models.


\subsection{What is a good req spec?}

 When is a requirements model ready? These are difficult questions that call for a definition of what we mean by a the quality of a requirements model. The quality of a requirements model is assessed during requirements validation and that activity seeks to answer if a requirements model is good enough or if we need to do more work in order to improve its quality.

Here are some examples of typical \textbf{quality factors} of requirements specifications that are assessed during validation:
\begin{itemize}
  \item \textbf{Correctness}. \TODO{one-liner def of each quality aspect}
  \item \textbf{Completeness}.
  \item \textbf{Conciseness}
  \item \textbf{Ambiguity}.
  \item \textbf{Verifiability}.
  \item \textbf{Consistency}
  \item \textbf{Traceability}
\end{itemize}



\subsection{Different kinds of requirements}

\begin{itemize}
  \item Explicit vs implicit requirements.
  \item The parts of information in entities may be combined/bundled/separated into the same or split across different in many ways. These kinds of requirements are actualy information aspects that can be part of the same requirement. One req can include both functions and quality. One req can include both data and logic. Etc. Combine abstration levels, levels of detail, levels of aggregation, different degrees of formality etc.
\end{itemize}


\subsubsection{Functional Requirements (FR)}
Functional requirements deals with the expected logical behavior and the related data that is needed by the system to do what it should.
\begin{itemize}
  \item \textbf{Data}: information stored by the system, including domain-specific data entities and the format of input and output data. Data requirements often involves a data base or some other form of persistent storage of domain-specific data. Data requirements are related to user interface design, as users expect to see relevant data when they interact with the system to should help them accomplish their tasks. Data requirements usually includes some abstract representation of the internal system state, at-least those parts of the system state that are externally observable.
  \item \textbf{Logic}: the functional relation between input and output. The modeling of requirements on the logical behavior of a system involves the analysis of what is the desired output of a specific computation given some expected input in a specific usage context, and also what happens if unexpected input is given. We want to specify logic at a higher abstraction level than the actual implementation in source code, while capturing the underlying intent of the requested functionality. 
\end{itemize}

Behavioral requirements are often called ''business logic'' or just ''business requriements''. Data requirements on domain-related information are often less complex and easier to elicit, specify and validate, compared to requirements that deal with intricate business logic in a multi-faceted usage context. Here we use the term ''functional requirements'' to cover both requirements on data and logical behavior, but sometimes data requirements are treated as a its own kind of requirements, in separation of (other) functional requirements. 

\subsubsection{Quality Requirements (QR)}

Aka NFR.

\begin{itemize}
  \item \textbf{Accuracy}. \TODO{on-liner def of each quality aspect}
  \item \textbf{Capacity}.
  \item \textbf{Performance}.
  \item \textbf{Reliability}.
  \item \textbf{Usability}.
  \item \textbf{Security}.
  \item \textbf{Safety}.
\end{itemize}

Quality and functionality are often intertwined and quality often manifest itself into behavior logic and data that support the achievement of a system quality. As an example consider system security. In order to have high security we may want to implement a user login function implying that the system needs to store data about user identities and secret passwords. As quality and functionality often come hand-in-hand, it is often wise to specify these aspects together by letting feature models include both functional and quality aspects. However, many quality aspects are cross-cutting, with system-wide impact, so we also often need a high-level, verifiable model of quality at system level. 

A further complication is that quality aspects often work against each other. For example, achieving high security often involved implementing features such as authentication and encryption, while these features can negatively impact the usability of the system.

\subsubsection{Requirements at different levels}

\begin{itemize}
  \item Levels of Abstraction: the Goal-Design scale
  \item Levels of Detail: amount and richness of information
  \item Levels of Aggregation: grouping and linking
  \item Levels of Formality: from unstructured to mathematical
  • Very informal: free-form
  representation, no explicit rules
  • Very formal: syntax, semantics,
  inference, meta-language
  • Formality enables automatic checks,
  concise models, ...
  • Formalization requires effort,
  knowledge, skills, ...
\end{itemize}


\subsection{Feature requirements}%

Textual features combining data and logic. Releasable unit of FR+QR.

The most common type of requirement.

Flexible and expressive but ambiguous.

Boosted by LLM:s

\subsection{Specification tools}

\subsubsection{Natural Language}

-- controlled domain-terminology 

-- explicit or implicit reqs, 

-- shall / should / must / optional MOSCOW EARS

-- point to RM chapter on textual feature requirements

\subsubsection{Generic tools}
  word processors, presentation tools, drawing tools , issue trackers 

  AI and Large language models 

\subsubsection{RE-specific tools}
  Features of RE Tools: Store, Collaborate, Document Model, Language, traceability, different views, RE-aware AI support, ...

  Examples: Doors, RequisitePro, reqT

\subsubsection{A requirements metalanguage}

The reqT meta-model cheat-sheet





%!TEX encoding = UTF-8 Unicode
%!TEX root = ../book/reqeng-book.tex

\chapter{Context}%

\section{Context of RE}

\subsection{Business model}
Funding model. Economy of dev.

License fee. Freemium. Subscription.

\subsection{Customer-supplier relationship}

\subsection{Product and integration}

Type of product

\begin{itemize}
  \item Level of customization (customer specific, generic, mixed)
  \item Hardware integration (Pure SW, embedded HW+SW, pure HW)
  \item Network integration (stand-alone, network connection, cloud: back-end -- front-end)
\end{itemize}

\subsection{Service and delivery}

\url{https://en.wikipedia.org/wiki/Product_lifecycle}

\url{https://en.wikipedia.org/wiki/Software_as_a_service}

Delivery model: local executable, cloud service, on-prem self-hosting, saas, etc etc


\section{Domain modeling}

\subsection{Stakeholder modeling}

\subsection{Goal modeling}

Helps or Hurts.

Related to both elicitation (stakeholder intentions) and specification (user roles and goals).

\TODO{Decide how much on Goals here or more details under Elicitation}???

\subsection{Context diagram}

Inner and outer domain.

\TODO{Decide if it should be called 'Context Diagram' instead?}






%!TEX encoding = UTF-8 Unicode
%!TEX root = ../book/reqeng-book.tex

\chapter{Elicitation}%
blabla


Elicitation Goals \& Challenges%

\section{What information to elicit?}

\begin{itemize}
  \item Present
  \item Future 
  \item Consensus 
  \item Commitment 
  \item Innovative ideas
  \item Goals -- Domain -- Product -- Design
  \item Priorities (cost, benefit, risk of requriements)
\end{itemize}

\section{Why is elicitation so hard?}

Elicitation challenges (barriers/difficulties) 

\section{Interacting with Stakeholders}%

\subsection{Surveys}

\subsection{Interviews}

blabla

\subsection{Case Studies}

Demonstrations

Observation of Current work 

Prototyping -> its own chapter 

Usability testing -> own subchapter in Module on Quality

Pilot system deployment

\subsection{Requirements from Operation}

support issues, telemetry, user forums and social media, A/B-testing experiments


\section{Business Intelligence}
blabla  (sv Omvärldsanalys)

study similar companies, competitors, subcontractors, suppliers 

\section{Creativity and Innovation}

trigger ideas 

assess novelty and market 

brainstorming

focus-groups

%!TEX encoding = UTF-8 Unicode
%!TEX root = ../book/reqeng-book.tex

\chapter{Prioritization}%
blabla

\section{Why prioritize?}%

\section{The prioritization process}

\subsection{Activities}
who does what when 

\subsection{Challenges}

Finding a good abstraction level

Combinatorial explosion

Inter-dependencies

Not easy to predict the future

Power and politics

\subsection{Aspects}

Example aspects:

-- Importance (e.g. financial benefit, urgency, strategic value, market share...)

--  Penalty (e.g. bad-will if requirement not included)

--  Cost (e.g., staff effort)

--  Time (e.g., lead time)

--  Risk (e.g., technical risk, business risk)

--  Volatility (e.g. scope instability, probability of change)

-- Other aspects: competitors, brand fitness, competence, release theme...

-- Combination of aspects: cost vs. benefit, cost vs. risk, importance vs. volatility

-- Optimize: minimize or maximize some combinations, e.g., cost vs. benefit

\section{Scales}

\section{Methods}

\textbullet Grouping, Numerical assignment (grading)

-- Can be done using any scale (categorical, ordinal, ratio)

-- Quick and easy; but a risk is that all reqs are deemed highly important as they are
not challenged against each other; may be misinterpreted as ratio scale (even if
”4” not necessarily is "twice as much" as ”2” when using an ordinal scale).

\textbullet 100-dollar-test

-- Ratio scale, quick and easy, risk of shrewd tactics (listigt taktikspel)

\textbullet Ranking (sorting)

-- Ordinal scale, pairwise comparison, easy and rather quick

\textbullet Top-ten (or Top-n)

-- Ordinal scale if the top list is ranked / Categorical grouping if not ranked; very
quick and simple, gives a rough estimate on a limited set of req

\textbullet (Analytical Hierarchy Process (AHP))

-- Ratio scale, pairwise comparison, tool is needed, redundancy gives measure of
consistency 


\part{Functionality}  %%%%%%%%%%%%%%%%%%%%%%%%%%%%


%!TEX encoding = UTF-8 Unicode
%!TEX root = ../book/reqeng-book.tex

\chapter{Data Requirements}%
blabla


\section{Data Dictionary}
blabla

\section{Entity-Relationship Diagrams}
blabla

\section{Class Diagrams}
blabla

\section{Data Screens}
blabla  (aka Virtual Windows)

\section{Formal Data Specification}%

\subsection{Regular Expressions}

\subsection{Protocol Buffers}

\url{https://en.wikipedia.org/wiki/Protocol_Buffers}

%!TEX encoding = UTF-8 Unicode
%!TEX root = ../book/reqeng-book.tex

\chapter{Logic}%

Functional Requirements on business logic that provide constraints on system behavior.


\section{When to decide user-system work split?}

\section{Contextual usage modeling}%

\subsection{User stories}

A user story is an informal, general explanation of a software feature written from the perspective of the end user. Its purpose is to articulate how a software feature will provide value to the customer.

\url{https://en.wikipedia.org/wiki/User_story}

\url{https://www.atlassian.com/agile/project-management/user-stories}

\subsection{Use cases and tasks}

\subsection{Narratives and personas}

Narrative: Rich and detailed description of a real-world usage situation in the form of a short-story.

Persona: a fictional character representing an imagined human user, often attributed with concrete personal attitudes and desires. 


\section{System behavior modeling}

-- Domain Events

-- Product Events

-- Product Functions triggered by Events: (Input, State) => Output

\subsection{State diagram}%
\url{https://en.wikipedia.org/wiki/State_diagram} 

\subsection{Interaction diagrams}

\begin{itemize}
  \item UML Sequence Diagrams \url{https://en.wikipedia.org/wiki/Sequence_diagram}
  \item ITU Message Sequence Charts \url{https://en.wikipedia.org/wiki/Message_sequence_chart}
\end{itemize}

\subsection{Data-flow diagrams}%
\url{https://en.wikipedia.org/wiki/Data-flow_diagram} 





\chapter{Prototyping}

\section{Why prototype?}

What is it?

Can support elicitation, but also specification, validation and selection. Explain how...

Why do it? Purpose of prototyping.

\begin{itemize}
  \item Exploration and learning
  \item Communication: sales, aligment
  \item Incremental development
  \item Quality improvement
  \item Validation and testing
  \item Problem-solution fit, product-market fit, technical feasibility, usability testing
\end{itemize}

\section{Prototype Scoping}

\begin{itemize}
  \item Breadth of functionality
  \item Functional refinement
  \item Visual appearance
  \item Interactive and haptic behavior
  \item Data realism
\end{itemize}


\section{Prototype Media}
paper or computer-based
\begin{itemize}
  \item Sketch: 
  \item Wireframe
  \item Mockup
  \item Source-code software
  \item Other: video, interview
\end{itemize}

\section{Prototype Usage}
\begin{itemize}
  \item Reviewers: internal, FFF, external
  \item Prototype interaction: yes no (demo)
  \item Review approach: scenario-based, free
  \item Usage environment
\end{itemize}

\section{Exploration Strategy}
\begin{itemize}
  \item Single vs parallell Exploration
  \item Iteration focus: business, product, feature, optimization
  \item Iteration size
\end{itemize}


\chapter{Delegated Requirements}

\section{Standards as requirements}

\section{Regulatory requirements}

\subsection{General Data Protection Regulation} 

\subsection{Artificial Intelligence Act}
Point to chapter on AI-act in "Special Domains"

\section{Test cases as requirements}

\section{Development process requirements}


\part{Quality} %%%%%%%%%%%%%%%%%%%%%%%%%%%%

\chapter{Quality Requirements}
\section{What is product quality?}
Different Quality Models. Relation between product-level quality and domain-level quality.

\url{https://en.wikipedia.org/wiki/ISO/IEC_9126}

\section{Product-level qualities}
\subsection{Accuracy}

Accuracy of data representation. Relates to Data model and  representation of data fields in records.

Reqs on: How to handle rounding errors?

\subsection{Capacity}
Subcategory of efficiency.

\subsection{Performance}
Subcategory of efficiency.

\subsection{Reliability}

\subsection{Maintainability}

\section{Domain-level qualities}

\subsection{Usability and UX}

Interaction requirements, quality of interaction design.

\subsection{Safety and Security}

\url{https://en.wikipedia.org/wiki/Information_security}

\chapter{Requirements Validation}

\section{Specification quality aspects}

\TODO{decide how much here and how much in intro on spec quality}

\section{Requirements inspections}
static validation of req models

capture-recapture

fault seeding

\section{Dynamic validation}

Dynamic validation methods:
\begin{itemize}
  \item Usability testing
  \item Prototyping
  \item Simulation
\end{itemize}

\subsection{Usability testing}

\subsection{Simulation}

\TODO{Or should Similuation get it's on chapter after Prototyping in Part II Functionality???}

\chapter{Product scoping}

\section{Release themes and epics}

bundles of features and use cases, a bundle that delivers value

\section{Requirements status ladder}

\section{Release planning} 

\chapter{Product verification}
\section{Unit testing}
\section{System testing}
\section{Acceptance testing}
\section{Regression testing}

\part{Special Contexts}

\chapter{Special Development Modes}

\section{Agile RE}

\section{Open Source RE}

\section{Continuous RE}

\subsection{CI/CD}
\subsection{DevOps}

\chapter{Special Product Types}

\section{RE for High-Assurance Systems}
\subsection{Information security audit}
\url{https://en.wikipedia.org/wiki/Information_security_audit}

\section{RE for Artificial Intelligens}

\section{RE for Embedded Systems}
\subsection{Hardware-Software co-design}

\part{Appendicies}\appendix

\chapter{Introduction to reqT}

\chapter{Lab 1: Prioritization}

\chapter{Lab 2: Release Planning}

\chapter{Course Project}

\section{Literature Studies}

\section{System Ideas}

\section{Course Project Outline}
\subsection{Management roles}
\subsection{Release 1: Scoped}
\subsection{Release 2: Specified}
\subsection{Release 3: Validated}
\subsection{Assessment Matrix}


\chapter{Learning from Public RE Disasters}

\section{Swedish National Police System}

\url{https://computersweden.se/article/1304245/haveriet-inifran-sa-gick-pust-fran-succe-till-fiasko.html}

\section{Swedish Regional Health Care System}

\url{https://computersweden.se/article/1276485/miljardprojekt-i-varden-underskattade-lagstiftning-och-forsenades-rejalt.html}

\url{https://lakartidningen.se/aktuellt/nyheter/2022/04/skane-staller-krav-pa-cerner-om-systemet-millennium/}

\url{https://lakartidningen.se/aktuellt/nyheter/2022/02/nytt-journalsystem-drar-ut-pa-tiden-vgr-vill-ha-pengar-tillbaka/}

\url{https://www.svt.se/nyheter/lokalt/vast/brister-i-nya-patientsystemet-har-kostat-skattebetalarna-en-miljard}

\end{document}
